\section{实验结果}

\subsection{候选池质量}

表~\ref{tab:enhanced-bioasq-main-results} 总结了当前 BioASQ 留出测试集主结果。原始 Hybrid RRF 测试基线的 Recall@10 为 0.4636，MRR@10 为 0.7550，nDCG@10 为 0.5848，Recall@100 为 0.6077。增强 w122 候选池根据验证集 Recall@100 选择，并在测试集达到 Recall@100 = 0.7388，为重排序提供了更强的候选上限。这一步很重要，因为如果金标准证据没有进入候选池，后续重排序无法恢复该证据。

\begin{table}[t]
\centering
\small
\caption{增强候选池上的 BioASQ 留出测试证据检索结果。}
\label{tab:enhanced-bioasq-main-results}
\begin{tabular}{lrrrr}
\toprule
方法 & Recall@10 & MRR@10 & nDCG@10 & Recall@100 \\
\midrule
原始 Hybrid RRF & 0.4636 & 0.7550 & 0.5848 & 0.6077 \\
增强 Hybrid w122 & 0.4660 & 0.7530 & 0.5844 & 0.7388 \\
MedCPT Cross-Encoder & 0.5172 & 0.7775 & 0.6390 & 0.7388 \\
增强 KCH-MedRank & 0.5332 & 0.7882 & 0.6446 & 0.7388 \\
\bottomrule
\end{tabular}
\end{table}



\subsection{留出测试集主结果}

Enhanced KCH-MedRank 达到 Recall@10 = 0.5332、MRR@10 = 0.7882、nDCG@10 = 0.6446、Recall@100 = 0.7388。相对于原始 Hybrid RRF，这对应 Recall@10 +0.0696、MRR@10 +0.0332、nDCG@10 +0.0598 的提升。相对于增强 Hybrid RRF 的候选排序，KCH-MedRank 的 Recall@10 提升 +0.0673，MRR@10 提升 +0.0352，nDCG@10 提升 +0.0602，evidence coverage@10 提升 +0.0673。

表~\ref{tab:enhanced-bioasq-bootstrap-hybrid} 报告了相对于增强 Hybrid RRF 的成对 bootstrap 检验。四个 top-10 指标均达到显著，$p=0.0002$。这说明提升并不只是来自更强候选池，重排序步骤本身也显著改变了 top-10 证据排序。

\begin{table}[t]
\centering
\small
\caption{BioASQ 留出测试集上增强 KCH-MedRank 与增强 Hybrid RRF 的成对 bootstrap 比较。}
\label{tab:enhanced-bioasq-bootstrap-hybrid}
\begin{tabular}{lrrrr}
\toprule
指标 & Hybrid RRF & KCH-MedRank & 95\% CI & p 值 \\
\midrule
Recall@10 & 0.4660 & 0.5329 & [+0.0548, +0.0800] & 0.0002 \\
MRR@10 & 0.7530 & 0.7867 & [+0.0207, +0.0470] & 0.0002 \\
nDCG@10 & 0.5844 & 0.6433 & [+0.0495, +0.0687] & 0.0002 \\
Evidence coverage@10 & 0.4660 & 0.5329 & [+0.0548, +0.0800] & 0.0002 \\
\bottomrule
\end{tabular}
\end{table}


\subsection{与生物医学语义重排序的比较}

在同一增强候选池上，真实 MedCPT Cross-Encoder 基线达到 Recall@10 = 0.5172、MRR@10 = 0.7775、nDCG@10 = 0.6390。表~\ref{tab:enhanced-bioasq-significance} 报告了相对于该强语义基线的成对 bootstrap 结果。KCH-MedRank 在 Recall@10 上显著提升（$p=0.0074$），MRR@10 和 nDCG@10 为正向提升但未达到统计显著。

\begin{table}[t]
\centering
\small
\caption{同一增强候选池上增强 KCH-MedRank 与 MedCPT Cross-Encoder 的成对 bootstrap 比较。}
\label{tab:enhanced-bioasq-significance}
\begin{tabular}{lrrrr}
\toprule
指标 & 基线 & KCH-MedRank & 差值 & p 值 \\
\midrule
Recall@10 & 0.5172 & 0.5329 & +0.0157 & 0.0076 \\
MRR@10 & 0.7775 & 0.7867 & +0.0093 & 0.2060 \\
nDCG@10 & 0.6390 & 0.6433 & +0.0043 & 0.3840 \\
\bottomrule
\end{tabular}
\end{table}


与 MedCPT Cross-Encoder 的比较是对所提出重排序器更严格的检验。结果表明，KCH-MedRank 能在强生物医学语义重排序器之外进一步提高 top-10 证据覆盖，但 MRR 和 nDCG 的排序敏感指标只能表述为正向趋势，而不是显著提升。这些结果支持本文的证据检索主张，但尚不足以把答案生成作为主要贡献。

\subsection{答案支持诊断}

PubMedQA 答案选择诊断进一步区分了证据检索与答案正确性。对于 Hybrid 和 HGB 的 top 10 证据，规则化引用支持率均为 1.0000，因为每个评估问题都至少检索到一个金标准证据段落。然而，答案支持率仍受轻量 yes/no/maybe 答案选择器限制：多数类基线的支持答案比例为 0.5550，而 TF-IDF 逻辑回归在 Hybrid 与 HGB top 10 上分别为 0.5050 和 0.4800。在 98 个抽取到问题实体的问题中，top 10 答案-证据实体一致性为 0.5510。这些诊断使用 qrels、金标准标签和确定性实体重叠，而不是 LLM 评审。因此，在加入更强答案与忠实性建模前，答案生成仍应作为次要任务。
